%%%%%%%%%%%%%%%%%%%%%%%%%%%%%%%%%%%%%%%%%%%%%%%%%%%%%%%%%%%%%%%%%%%%%%%
%                                                                     %
%              Scientific Word   Wrap/Unwrap  Version 2.5             %
%                                                                     %
% If you are separating the files in this message by hand, you will   %
% need to identify the file type and place it in the appropriate      %
% directory.  The possible types are: Document, DocAssoc, Other,      %
% Macro, Style, Graphic, PastedPict, and PlotPict. Extract files      %
% tagged as Document, DocAssoc, or Other into your TeX source file    %
% directory.  Macro files go into your TeX macros directory. Style    %
% files are used by Scientific Word and do not need to be extracted.  %
% Graphic, PastedPict, and PlotPict files should be placed in a       %
% graphics directory.                                                 %
%                                                                     %
% Graphic files need to be converted from the text format (this is    %
% done for e-mail compatability) to the original 8-bit binary format. %
%                                                                     %
%%%%%%%%%%%%%%%%%%%%%%%%%%%%%%%%%%%%%%%%%%%%%%%%%%%%%%%%%%%%%%%%%%%%%%%
%                                                                     %
% Files included:                                                     %
%                                                                     %
% "/document/chap19.tex", Document, 12432, 4/14/1999, 23:02:56, ""    %
%                                                                     %
%%%%%%%%%%%%%%%%%%%%%%%%%%%%%%%%%%%%%%%%%%%%%%%%%%%%%%%%%%%%%%%%%%%%%%%

%%%%%%%%%%%%%%%%%%%%%% Start /document/chap19.tex %%%%%%%%%%%%%%%%%%%%%

%% This document created by Scientific Notebook (R) Version 3.0


\documentclass[12pt,thmsa]{article}
%%%%%%%%%%%%%%%%%%%%%%%%%%%%%%%%%%%%%%%%%%%%%%%%%%%%%%%%%%%%%%%%%%%%%%%%%%%%%%%%%%%%%%%%%%%%%%%%%%%%%%%%%%%%%%%%%%%%%%%%%%%%
\usepackage{sw20jart}

%TCIDATA{TCIstyle=article/art4.lat,jart,sw20jart}

%TCIDATA{<META NAME="GraphicsSave" CONTENT="32">}
%TCIDATA{Created=Mon Aug 19 14:52:24 1996}
%TCIDATA{LastRevised=Wed Apr 14 16:02:55 1999}
%TCIDATA{CSTFile=Lab Report.cst}
%TCIDATA{PageSetup=72,72,72,72,0}
%TCIDATA{AllPages=
%F=36,\PARA{038<p type="texpara" tag="Body Text" >\hfill \thepage}
%}


\input{tcilatex}
\begin{document}


\vspace{1pt}Lagrange Multipliers

Chapter 19 \ Chong/Zak

Author: Robert Beezer

History:

1999/04/07\qquad First version.

\vspace{1pt}

\subsubsection{Problems}

\vspace{1pt}

\begin{enumerate}
\item  min $f(x)$, subject to \ $h_{i}(x)=0$, \ $1\leq i\leq m$, \ $%
g_{j}(x)=0$, \ $1\leq j\leq p$. \ \ We can ``roll-up'' the \ $h$'s and the \
$g$'s into vector-valued functions and compare them to the zero vector. \ \
\ Note that this problem formulation encompasses a wide variety of problems,
including linear programming. \ At first we'll stick with the equality
constraints only.\bigskip

\item  Example (Edwards and Penney, 4e, Sec 14.9):

min $\ \ x^{2}+y^{2},\ $subject to $\ \ xy=1$.

Solve graphically with circles as level curves, touching up to a rotated
hyperbola at points like $(1,1)$ \ and \ $(-1,-1)$.\bigskip

\item  Example (Edwards and Penney, 4e, Example 14.9.4):

Three variables, two equality constraints.

The plane \ $x+y+z=12$ \ intersects \ $z=x^{2}+y^{2}$ \ to form an ellipse.
\ Where is the lowest point on this intersection?

min $z$, \ subject to \ \ $(x+y+z-12,$\ $x^{2}+y^{2}-z)=(0,0)$. \ Draw a
picture - $(2,2,8)$ \ is the low point.\bigskip

\item  Example (Uhl/Peressini/Sullivan, Problem 7.4):

Three variables, two equality constraints.

min \ $\frac{1}{x^{2}+y^{2}+z^{2}}$, \ subject to $x^{2}+2y^{2}+3z^{2}=1$, \
$x+y+z=0$.

\emph{Discuss} the appropriate picture.
\end{enumerate}

\vspace{1pt}

\subsubsection{Theory}

\vspace{1pt}

\begin{enumerate}
\item  Definition: \ A point \ $x^{\ast }$ is \emph{feasible} if \ $%
h(x^{\ast })=0$.\bigskip

\item  Definition: \ Given \ $h(x)=(h_{1}(x),h_{2}(x),\ldots h_{m}(x))$ \
where \ $x$ \ is an \ $n$-slot vector, the Jacobian is the \ $m\times n$ \
matrix whose rows are the gradients \ $\nabla h_{i}(x)$. \ The notation is \
$Dh(x)$. \ \ Illustrate with last example above \ (U/P/S 7.4).
\[
Dh(x)=\left(
\begin{array}{lll}
2x & 4y & 6z \\
1 & 1 & 1
\end{array}
\right)
\]
\bigskip

\item  Definition: \ A feasible point \ $x^{\ast }$ is \emph{regular} if the
set of gradients \ $\{$\ $\nabla h_{i}(x)\mid $\ $1\leq i\leq m\}$ \ is
linearly independent. \ \ (Which would imply that \ $m<n$, so we often
assume this is the case). \ This says that the Jacobian is a matrix of full
rank.\bigskip

\item  Definition: \ The tangent space at a point \ $x^{\ast }$ \ on the
surface \ $\{x\mid h(x)=0\}$ \ \ is the set
\[
T(x^{\ast })=\{y\mid Dh(x^{\ast })y=0\}
\]

the null space of the Jacobian.\bigskip

\item  Example: \ (last example above U/P/S 7.4) at \ $x^{\ast
}=(5/10,-6/10,1/10)$

Check feasibility: \ EZ

Check regularity: $\ \ \ Dh((5/10,-6/10,1/10))=\left(
\begin{array}{lll}
1 & \frac{-24}{10} & \frac{6}{10} \\
1 & 1 & 1
\end{array}
\right) $. \ The rows are linearly indpendent. \ The only time they wouldn't
be is if \ \ $x^{\ast }=\alpha (\frac{1}{2},\frac{1}{4},\frac{1}{6})$. \ \
The linear constraint would then force \ $\alpha =0$, but then \ $x^{\ast }=0
$ \ and the quadratic constraint is not satisfied. \ This means that every
feasible point is also regular.

Check the tangent space:$\left(
\begin{array}{lll}
1 & \frac{-24}{10} & \frac{6}{10} \\
1 & 1 & 1
\end{array}
\right) \bigskip $, nullspace basis: $\left\{ \left(
\begin{array}{c}
\frac{15}{2} \\
1 \\
-\frac{17}{2}
\end{array}
\right) \right\} $

\item  Definition: \ the tangent plane is formed by shifting the tangent
space to \ $x^{\ast }$
\[
TP(x^{\ast })=x^{\ast }+T(x^{\ast })=\{x^{\ast }+y\mid y\in T(x^{\ast })\}
\]
\bigskip

\item  In the prior example this would be:$(5/10,-6/10,-1/10)+\alpha \left(
\begin{array}{ccc}
\frac{15}{2} & 1 & -\frac{17}{2}
\end{array}
\right) $, transpose:  \ for \ all \ $\alpha $. \ Visualize this as a line
perpendicular to the set at \ \ $x^{\ast }$.\bigskip

\item  Definition: \ The normal space at \ $x^{\ast }$ \ is:
\[
N(\ x^{\ast })=\{Dh(x^{\ast })^{T}z\mid z\in R^{m}\}
\]

Note carefully that this involves the transpose of the Jacobian, and is the
range of this matrix. \ It is all linear combinations of the rows of the
Jacobian, which are gradients of the constraint vector. \ Again, we can
shift this to the point \ $\ x^{\ast }$ \ to form what is called the normal
plane. \ Illustrate this some with the above example, writing linear
combinations of the rows of the Jacobian, translated to $\ x^{\ast }$.

\item  Fact: \ If \ $x\in N(\ x^{\ast })$ \ and \ $y\in T(\ x^{\ast })$ \ \
then \ they are orthogonal vectors. \ So these \emph{sets} are orthogonal -
each is precisely the set of vectors perpendicular to everything in the
other (which is not exactly what is shown below).

Proof:
\[
y^{T}x=y^{T}Dh(x^{\ast })^{T}z=(Dh(x^{\ast })y)^{T}z=0z=0
\]

\item  Fact: $\ \ N(\ x^{\ast })\cap T(\ x^{\ast })=\{0\}$

Proof: \ Presume \ $x$ \ is in the intersection, then by viewing it as being
in one or the other of the two sets, the above result shows that \ $xx=0$,
implying that \ $x$ \ is the zero vector.\bigskip

\item  Check the dimensions of these two spaces, and consider a basis for
each. \ Together they must span all of \ $R^{n}$. \ So any vector at all can
be written (uniquely) as a vector from $N(\ x^{\ast })$ \ added to a vector
from \ $T(\ x^{\ast })$.
\end{enumerate}

\vspace{1pt}

\subsubsection{Lagrange's Theorem}

\vspace{1pt}

\begin{enumerate}
\item  Read Theorem 19.2 (and the \emph{preceding} text) for a prrof in the
2-D, one constraint case.\bigskip

\item  Theorem: \ Suppose that \ $x^{\ast }$ \ is a regular point and a
local minimizer for \ $f(x)$ \ when subject to the constraint that \ $h(x)=0$%
. \ Then there exists a vector\ $\lambda ^{\ast }$ \ such that
\[
Df(x^{\ast })+\lambda ^{\ast }Dh(x^{\ast })=0
\]

\item  Notes: \ (1) \ Check the sizes and types of the objects above. \ \
(2) \ This is a necessary condition, but not sufficient. \ It will narrow
our search for minimizers, but not say anything about if these points are
definitely minimizers (some will be maximizers, some will be neither).

\item  Proof: \ Choose any tangent vector to the constraint surface: \ $y\in
T(x^{\ast })$. \ Theorem 19.1 (which we have not discussed) provides a
parametrized curve on the constraint surface that passes through \ $x^{\ast }
$ \ and has a derivative equal to the tangent there. \ That is, there exists
\ $x(t)$, \ \ differentiable curve such that \ \ $h(x(t))=0$ \ and at some
point \ $t^{\ast }$ \ we have \
\[
x(t^{\ast })=x^{\ast }\text{, \ \ \ \ }\left( \frac{dx}{dt}\right)
_{t=t^{\ast }}=y
\]

Define \ $\phi (t)=f(x(t))$. \ \ \ If \ $x^{\ast }$ \ is a local minimizer
of \ $f(x)$, \ then \ $t^{\ast }$ \ is a local minimizer of \ $\phi (t)$. \
\ Then
\[
\frac{d}{dt}\phi (t)=Df(x(t))\cdot \frac{dx}{dt}
\]

Evaluating at \ $x=t^{\ast }$ \ and using the FONC for a minimizer yields
\[
0=\frac{d}{dt}\phi (t^{\ast })=Df(x(t^{\ast }))\cdot \left( \frac{dx}{dt}%
\right) _{t=t^{\ast }}=\nabla f(x(t^{\ast }))\cdot y=\nabla f(x^{\ast
})\cdot y
\]

So any \ $y\in T(x^{\ast })$ \ satisfies \ $\nabla f(x^{\ast })\cdot y=0$.

So \ $\nabla f(x^{\ast })$ \ is an element of the set orthogonal to $%
T(x^{\ast })$, \ namely $N(\ x^{\ast })$. \ \ So \ $\nabla f(x^{\ast
})=Dh(x^{\ast })^{T}z$ \ for some \ $z$. \ Rearranged this becomes
\[
\nabla f(x^{\ast })-Dh(x^{\ast })^{T}z=0
\]

so the vector described by the statement of the theorem is \ $-z$, \ that is
\ $\lambda ^{\ast }=-z$.
\end{enumerate}

\vspace{1pt}

\vspace{1pt}Problems Again

\vspace{1pt}

\begin{enumerate}
\item  Example (Edwards and Penney, 4e, Sec 14.9):

min $\ \ x^{2}+y^{2},\ $subject to $\ \ xy=1$.
\[
(2x,2y)-\lambda (y,x)=(0,0)\Rightarrow \frac{\lambda }{2}=\frac{y}{x}=\frac{x%
}{y}
\]

So
\[
x^{2}=y^{2}\text{ \ and thus \ }x=y=\pm 1
\]

Discuss why these two points, \ $(1,1)$ \ and \ $(-1,-1)$, \ must be
minimizers.\bigskip

\item  Example (Edwards and Penney, 4e, Example 14.9.4):

The plane \ $x+y+z=12$ \ intersects \ $z=x^{2}+y^{2}$ \ to form an ellipse.
\ Where is the lowest point on this intersection?

min $z$, \ subject to \ \ $(x+y+z-12,$\ $x^{2}+y^{2}-z)=(0,0)$.
\[
\left(
\begin{array}{r}
0 \\
0 \\
1
\end{array}
\right) -\left(
\begin{array}{rr}
1 & 2x \\
1 & 2y \\
1 & -2z
\end{array}
\right) \left(
\begin{array}{r}
\lambda _{1} \\
\lambda _{2}
\end{array}
\right) =\left(
\begin{array}{r}
0 \\
0 \\
0
\end{array}
\right)
\]

First two slots together yield that
\[
x=y=\frac{\lambda _{1}}{2\lambda _{2}}
\]

This is enough to convert the two constraints to
\[
2x+z=12\text{ \ and \ }z=2x^{2}
\]

which together yield \ $x=2$ \ and \ $x=-3$. \ These two solutions lead
(respectiveley) to \ $z=8$ \ and \ $z=9$, ostensibly a max and a min. \
These solution techniques are typical - try not to determine the $\lambda $%
's if you can avoid it, work on the variables of the problem.\bigskip

\item  Example (Uhl/Peressini/Sullivan, Problem 7.4):

Three variables, two equality constraints.

min \ $\frac{1}{x^{2}+y^{2}+z^{2}}$, \ subject to $x^{2}+2y^{2}+3z^{2}=1$, \
$x+y+z=0$.
\[
\left(
\begin{array}{r}
\frac{-2x}{\left( x^{2}+y^{2}+z^{2}\right) ^{2}} \\
\frac{-2y}{\left( x^{2}+y^{2}+z^{2}\right) ^{2}} \\
\frac{-2z}{\left( x^{2}+y^{2}+z^{2}\right) ^{2}}
\end{array}
\right) -\left(
\begin{array}{rr}
2x & 1 \\
4y & 1 \\
6z & 1
\end{array}
\right) \left(
\begin{array}{r}
\lambda _{1} \\
\lambda _{2}
\end{array}
\right) =\left(
\begin{array}{r}
0 \\
0 \\
0
\end{array}
\right)
\]

Multiply through by $\left( x^{2}+y^{2}+z^{2}\right) ^{2}$ \ which is like
doing a maximum problem on the reciprocal of the objective function. \ The
equations that result are
\[
-2x-2x\lambda _{1}=\lambda _{2}\text{ \ and \ }-2y-4y\lambda _{1}=\lambda
_{2}\text{\ \ and }-2z-6z\lambda _{1}=\lambda _{2}
\]

Setting these expressions for $\lambda _{2}$ \ equal to each other yields
\[
y=\frac{1+\lambda _{1}}{1+2\lambda _{1}}x\text{ \ \ and \ }z=\frac{1+\lambda
_{1}}{1+3\lambda _{1}}x
\]

Substituting into \ $x+y+z=0$ \ creates the equation
\[
\left( 1+\frac{1+\lambda _{1}}{1+2\lambda _{1}}+\frac{1+\lambda _{1}}{%
1+3\lambda _{1}}\right) x=0\Rightarrow \left( \frac{3+12\lambda
_{1}+11\lambda _{1}^{2}}{\left( 1+2\lambda _{1}\right) \left( 1+3\lambda
_{1}\right) }\right) =0\text{ \ or \ }x=0
\]

Two solutions are \ \ $\lambda _{1}=\frac{-6\pm \sqrt{3}}{11}$. \ These can
be plugged into the expressions above for \ $y$ \ and \ $z$, and then into
the quadratic constraint \ to get a single equation in \ $x$. \ With \ \ $x$
\ and \ $\lambda _{1}$ \ we can determine \ \ $y$ \ and \ $z$, and evaluate
the function. \ \ These computations, from Mathematica, follow:

\begin{tabular}{|l|l|l|l|l|}
\hline
$\lambda _{1}$ & $x$ & $y$ & $z$ & $\frac{1}{x^{2}+y^{2}+z^{2}}$ \\ \hline
$\frac{-6+\sqrt{3}}{11}$ & $0.131633$ & $0.359627$ & $-0.49126$ & $2.57735$
\\ \hline
$\frac{-6+\sqrt{3}}{11}$ & $-0.131633$ & $-0.359627$ & $0.49126$ & $2.57735$
\\ \hline
$\frac{-6-\sqrt{3}}{11}$ & $0.661225$ & $-0.48405$ & $-0.177175$ & $1.42265$
\\ \hline
$\frac{-6-\sqrt{3}}{11}$ & $-0.661225$ & $0.48405$ & $0.177175$ & $1.42265$
\\ \hline
\end{tabular}

The actual minimum value is $\frac{11(746+425\sqrt{3})}{(5751+3296\sqrt{3})}%
=\allowbreak 1.\,\allowbreak 422\,649\,731$

In the case where \ $x=0$ \ we get all zeros, and an undefined functional
value.
\end{enumerate}

\vspace{1pt}

\vspace{1pt}

\vspace{1pt}

Problems: \ \ Chong/Zak \ \ Page 349 \ \ 1ab, 2a, 6

\end{document}

%%%%%%%%%%%%%%%%%%%%%%% End /document/chap19.tex %%%%%%%%%%%%%%%%%%%%%%